% TOP_PROBLEM: {"id": 64, "title": "Smale's 12th Problem: Centralizers of Diffeomorphisms", "category_id": 6, "category": {"id": 6, "name": "geometry", "display_name": "Geometry", "description": "Euclidean and non-Euclidean geometry, geometric structures.", "slug": "geometry", "order_index": 6, "created_at": "2026-07-31T15:26:25.671Z"}, "status": "open"}
% FIRST_POSTED: null
% ENTRY_KIND: statement_only
% Imported from: batch01.tex; UnsolvedMath ID 64
\documentclass[11pt]{article}
\usepackage[margin=25mm]{geometry}
\usepackage{fontspec}
\setmainfont{Latin Modern Roman}
\usepackage{amsmath,amssymb,mathtools}
\usepackage{unicode-math}
\setmathfont{Latin Modern Math}
\usepackage{xurl,hyperref}
\hypersetup{colorlinks=true,urlcolor=blue}
\setcounter{secnumdepth}{0}
\setlength{\parindent}{0pt}
\setlength{\parskip}{5pt}
\setlength{\emergencystretch}{3em}
\newcommand{\sref}[2]{\hyperlink{source:#1:#2}{[S#2]}}
\newcommand{\cataloguescope}[1]{\paragraph{Recorded target.}#1}
\begin{document}
% UnsolvedMath ID 64; source catalogue code SMA-012
\setcounter{section}{63}
\section{Smale's 12th Problem: Centralizers of Diffeomorphisms}\label{64}
{\small\sffamily UnsolvedMath ID 64 · Geometry}
\subsection{Definitions and mathematical statement}

\paragraph{Shared notation.}$\mathbb N=\{1,2,\ldots\}$, and $\mathbb Z,\mathbb Q,\mathbb R,\mathbb C$ denote the integers, rationals, reals and complex numbers. Any other conventions are specified locally.


Let $M$ be a closed connected smooth manifold, where closed means compact without boundary. Fix a differentiability class $r\in\mathbb N\cup\{\infty\}$. Let $\operatorname{Diff}^r(M)$ denote the group of bijective $C^r$ maps $M\to M$ whose inverses are also $C^r$, with its $C^r$ topology. For $f\in\operatorname{Diff}^r(M)$ its centralizer in the same group is
\[
 Z^r(f)=\{g\in\operatorname{Diff}^r(M):g\circ f=f\circ g\}.
\]
It always contains $\langle f\rangle=\{f^n:n\in\mathbb Z\}$, using inverse iterates when $n<0$. Here ``trivial centralizer'' means $Z^r(f)=\langle f\rangle$, not that $Z^r(f)$ consists only of the identity. Put
\[
 \mathcal T^r(M)=\{f\in\operatorname{Diff}^r(M):Z^r(f)=\langle f\rangle\}.
\]
A subset of this function space is residual if it contains a countable intersection of open dense subsets; a dense $G_\delta$ is such an intersection. \sref{64}{2}
\paragraph{Statement recorded in the supplied source.}
Determine whether, for each $M$ and $r$ as above, the diffeomorphisms with trivial centralizer have the following properties:
\begin{enumerate}
\item $\mathcal T^r(M)$ is dense in $\operatorname{Diff}^r(M)$;
\item $\mathcal T^r(M)$ is residual in $\operatorname{Diff}^r(M)$;
\item $\mathcal T^r(M)$ contains an open dense subset of $\operatorname{Diff}^r(M)$.
\end{enumerate}
These are the distinct genericity formulations of Smale's centralizer question, with the entire commuting group required to consist of iterates of $f$. \sref{64}{2}
For $r=1$, Source~S2 proves the residual assertion and, for positive-dimensional manifolds, recalls a negative answer to the open-dense assertion; these facts do not decide every higher-regularity case. \sref{64}{2}
\subsection{Short English statement}
Do diffeomorphisms whose only commuting diffeomorphisms are their own iterates form a dense set, a residual set, or a set containing an open dense subset, in each differentiability class?
\subsection{Sources}
\begin{description}
\item[\hypertarget{source:64:1}{S1}] ULAM.AI and the UnsolvedMath Contributors, UnsolvedMath: A Curated Collection of Open Mathematics Problems, record 64 (SMA-012). CC BY 4.0. \url{https://huggingface.co/datasets/ulamai/UnsolvedMath}.
\item[\hypertarget{source:64:2}{S2}] Christian Bonatti, Sylvain Crovisier and Amie Wilkinson, The C1 Generic Diffeomorphism Has Trivial Centralizer, Publications Mathématiques de l’IHÉS 109 (2009), 185--244; arXiv:0804.1416, Section 1.1, Question 1 and the following C1 results, pages 3--4. \url{https://arxiv.org/pdf/0804.1416}.
\end{description}
\label{end:64}
\end{document}
